\section{Выводы}

В ходе выполнения третьей части работы исследованы компиляционные методы
синтеза речевого сигнала и получены следующие результаты.

\begin{enumerate}
	\item Компиляционные методы синтеза речи строят сигнал из отрезков
	естественной речевой волны (фонем, аллофонов, дифонов). Для
	преобразования орфографического текста в фонемный используются правила
	транскрибирования: разложение букв Ю, Ё, Я, Е, редукция безударных
	гласных, оглушение и озвончение согласных, уподобление и смягчение.

	\item По четырём эталонным фразам (одна синтагма, перечисление, вопрос
	и утверждение) сформирована фонетическая база из 13 фонем. Показано,
	что изолированные гласные, записанные в виде списка, имеют сильно
	различающийся тон (119--205~Гц), что впоследствии порождает искажения
	при компиляции.

	\item Компиляция фраз из отдельных фонем даёт разборчивую, но
	неестественную речь: на стыках фонем возникают скачки $F_0$ до
	35--100~Гц, фразы монотонны (вопрос не отличим от утверждения), а
	коартикуляция между звуками отсутствует. В этом заключаются основные
	недостатки компиляционного метода при использовании слишком мелких
	единиц.

	\item Правка осциллограммы (замена неудачных гласных ядрами из слогов,
	сглаживание стыков, выравнивание уровней) уменьшила скачок $F_0$ во
	фразе <<Ты дома.>> с 35 до 15~Гц и приблизила тон к эталонному.

	\item Совершенствование фонетической базы --- переход к целым слогам
	--- сохранило естественные переходы <<согласный--гласный>> и снизило
	максимальный скачок $F_0$ в 3--8 раз (до 6--13~Гц), а также
	приблизило длительность и уровень фраз к эталонным.

	\item Просодическая обработка контурами $F_0(t)$, $A(t)$, $T(t)$ с
	сохранением формант придала фразам интонацию: вопрос <<Ты дома?>>
	получил восходящий контур ($+6{,}8$~полутона), утверждение ---
	нисходящий ($-1{,}5$~полутона), что сделало их различимыми на слух, а
	длительности приблизились к эталонным.

	\item Для получения естественной речи необходимо учитывать не только
	фонемный состав, но и просодию: непрерывность мелодического контура,
	ритм и энергетику. Чтобы <<синтезатор>> стал полноценной системой с
	неограниченным словарём, ему недостаёт анализа текста, автоматического
	транскриптора, просодической разметки, развитой базы аллофонов и
	дифонов, автоматического выбора единиц и сглаживания стыков.
\end{enumerate}

\begin{thebibliography}{4}
	\bibitem{lobanov}
		Лобанов, Б.~М. Речевой интерфейс интеллектуальных систем :
		учеб. пособие / Б.~М.~Лобанов, О.~Е.~Елисеева ; под науч. ред.
		В.~В.~Голенкова. --- Минск : БГУИР, 2006. --- 152 с.
	\bibitem{rylov}
		Рылов, А.~С. Анализ речи в распознающих системах /
		А.~С.~Рылов. --- Минск : Бестпринт, 2003. --- 264 с.
	\bibitem{rabiner}
		Рабинер, Л.~Р. Цифровая обработка речевых сигналов /
		Л.~Р.~Рабинер, Р.~В.~Шафер ; пер. с англ. --- М. : Радио и
		связь, 1981. --- 495 с.
	\bibitem{fant}
		Фант, Г. Акустическая теория речеобразования / Г.~Фант ; пер. с
		англ.; под ред. В.~С.~Григорьева. --- М. : Наука, 1964. --- 284 с.
	\bibitem{cloning}
		Лобанов, Б.~М. Компьютерное <<клонирование>> персонального голоса
		и речи / Б.~М.~Лобанов // Новости искусственного интеллекта. ---
		2002. --- №~5. --- С.~35--39.
	\bibitem{pecheneva}
		Виды языкового разбора : краткий комплексный
		словарь-справочник / Авт.-сост. Т.~А.~Печенева. --- Минск :
		Изд. центр БГУ, 2004. --- 99 с.
	\bibitem{zinder}
		Зиндер, Л.~Р. Общая фонетика / Л.~Р.~Зиндер. --- М. : Высш.
		шк., 1979. --- 312 с.
\end{thebibliography}
